\documentclass[french]{book}
\usepackage{teach}
\usepackage{verbatim}
\usepackage{amsmath,amsfonts,amssymb}
\usepackage{pgfplots}
\usepackage{mwe}
\RequirePackage{graphicx}
\RequirePackage{ifpdf}
 \ifpdf
   \DeclareGraphicsRule{*}{mps}{*}{}
 \fi


  \usepackage{fancyhdr}

  \usepackage{amsmath}
  \usepackage{amssymb}
  \usepackage{amsfonts}
  \usepackage{amsthm}
  \usepackage{mathrsfs}

  \usepackage{array}
  \usepackage{multicol}
  \setlength{\multicolsep}{3pt}

  \usepackage{graphicx}
  \graphicspath{{Figures/}}
  \usepackage{pstricks,pst-plot,pst-text,pst-tree,pst-eucl}
  \usepackage[all]{xy}

  \usepackage{fancybox}
  \usepackage{color}

%  \usepackage[frenchb]{babel}
  \usepackage{hyperref}

\usepackage{subfig}
\pgfplotsset{compat=newest}
\usepackage[all]{xy}
%\usepackage{geometry}
%\geometry{top=1cm, bottom=1cm, left=2cm, right=2cm}
\usepackage[utf8]{inputenc}
\usepackage[french]{babel}
\redefinecolor{thm}{title}{bgcolor}{alizarin}
\redefinecolor{prop}{title}{bgcolor}{meatbrown}
\redefinecolor{defn}{title}{bgcolor}{olivine}
\definecolor{alizarin}{rgb}{0.82, 0.1, 0.26}
\definecolor{meatbrown}{rgb}{0.9, 0.72, 0.23}
\definecolor{olivine}{rgb}{0.6, 0.73, 0.45}
\usepackage{mathrsfs}
%%
\newcommand{\R}{\mathbb {R}}
%\newcommand{\C}{\mathds {C}}
\newcommand{\N}{\mathbb {N}}
\newcommand{\Z}{\mathbb {Z}}
\newcommand{\Q}{\mathbb {Q}}
\newcommand{\D}{\mathbb {D}}
%%
\newcommand{\se}{\geq}
\newcommand{\ie}{\leq}
\newcommand{\tm}{\times}
%\newcommand{\ell}{l}
%\newcommand{\ell'}{l'}
%%
\newcommand{\ds}{\displaystyle}
\newcommand{\limn}{\lim_{n\rightarrow +\infty}}
\newcommand{\limo}[1][x]{\ds\lim_{#1\rightarrow 0}}
\newcommand{\limite}[2][x]{\ds\lim_{#1\rightarrow #2}}
\newcommand{\limpinf}[1][x]{\ds\lim_{#1\rightarrow +\infty}}
\newcommand{\limminf}[1][x]{\ds\lim_{#1\rightarrow -\infty}}
\newcommand{\limiteg}[2][x]{\ds\lim_{#1\xrightarrow{<}#2}}
\newcommand{\limited}[2][x]{\ds\lim_{#1\xrightarrow{>}#2}}
%%
\newcommand{\vect}[1]{\overrightarrow{#1}}
\newcommand{\oij}{(\textit{O},{\overrightarrow{i}},{\overrightarrow{j}})}
%
\newcommand{\cp}[2]{%
        \begin{pmatrix}
        #1\\
        #2
        \end{pmatrix}%
        }
%
\newcommand{\calb}{\mathscr{B}}
\newcommand{\calc}{\mathscr{C}}
\newcommand{\cald}{\mathscr{D}}
\newcommand{\cale}{\mathscr{E}}
\newcommand{\calf}{\mathscr{F}}
\newcommand{\calp}{\mathscr{P}}
\newcommand{\cals}{\mathscr{S}}
\newcommand{\calt}{\mathscr{T}}
%
\newcommand{\suite}[1][u]{\left( #1_{n} \right)_{n\in\N}}
\newcommand{\suitep}[1][u]{\left( #1_{n} \right)_{n\in\N^{*}}}
\usetikzlibrary{arrows}

\newcommand{\poly}[1][a]{#1_0+#1_1x+\cdots+#1_nx^n}


\begin{document}
\setcounter{chapter}{2}
\chapter{Dérivation des fonctions}
\section{Généralités}
\subsection{Limite en 0}
\begin{defn}
Soit $f$ une fonction définie sur un intervalle $I$ contenant 0 et $L \in \R$.

On dit que $f(x)$ tend vers $L$ quand $x$ tend vers 0 si on peut rendre $f(x)$ aussi proche de $L$ que l'on veut pour $x$ suffisamment proche de zéro. On note : $\limo[h] f(x)=L$
\end{defn}

\begin{exemple}
$\ds \limo[x]x^{2}=0$ \hspace{1em} $\ds \limo[x]x+1=1$ \hspace{1em} $\ds \limo[x]\sqrt{x}+x-2=-2$...

Déterminer la limite en 0 de $\dfrac{x^{2}-2x}{3x}$.
\end{exemple}


Pour tout $x\neq 0$, $\dfrac{x^{2}-2x}{3x}=\dfrac{x-2}{3}$ donc $\ds \limo[x]\dfrac{x^{2}-2x}{3x}=-\dfrac{2}{3}$.

%
\subsection{Nombre dérivé}
\begin{defn}
Soit $f$ une fonction définie sur un intervalle $I$. Soit $a \in I$ et $h \neq 0$ tel que $a+h \in I$.

On appelle taux de variation de $f$ entre $a$ et $h$ le réel $\dfrac{f(a+h)-f(a)}{h}$.

On dit que $f$ est dérivable en $a$ si, lorsque $h$ tend vers 0, le taux de variation de $f$ entre $a$ et $h$ tend vers un réel $L$ autrement dit si $\ds \lim_{h\rightarrow 0}\dfrac{f(a+h)-f(a)}{h}=L$.

Ce réel $L$ est appelé nombre dérivé de $f$ en $a$ et se note $f'(a)$.
\end{defn}

\begin{exemple}
Démontrer que la fonction $f : x \longmapsto x^{2}$ est dérivable en 2 et calculer $f'(2)$.

Démontrer que la fonction $g : x \longmapsto \sqrt{x}$ n'est pas dérivable en 0.
\end{exemple}


Pour tout $h\neq 0$, $\dfrac{f(2+h)-f(2)}{h}=\dfrac{(2+h)^{2}-4}{h}=\dfrac{h^{2}+4h}{h}=h+4$

Ainsi $\ds \limo[h]\dfrac{f(2+h)-f(2)}{h}=4$.

$f$ est donc dérivable en 2 et $f'(2)=4$.

Pour tout $h\neq 0$, $\dfrac{g(0+h)-g(0)}{h}=\dfrac{\sqrt{h}}{h}=\dfrac{1}{\sqrt{h}}$.

Or, $\dfrac{1}{\sqrt{h}}$ n'a pas de limite finie lorsque $h$ tend vers 0 donc $g$ n'est pas dérivable en 0.


\subsection{Interprétation graphique}
\begin{prop}
Soit $f$ une fonction définie sur $I$ et $C_{f}$ sa représentation graphique dans un repère. Soit $a\in I$.

Si $f$ est dérivable en $a$ alors $f'(a)$ est le coefficient directeur de la tangente à $C_{f}$ au point
$A(a;f(a))$. L'équation de cette tangente est alors : $y=f'(a)(x-a)+f(a)$.
\end{prop}

\begin{minipage}{9cm}
\underline{Illustration :} soit $A(a;f(a))\in C_{f}$.\par Soit $h \neq 0$ et $M(a+h;f(a+h))\in C_{f}$.

Le quotient $\dfrac{f(a+h)-f(a)}{h}$ représente le coefficient directeur de la droite $(AM)$. Lorsque $h$ tend vers 0, $M$ se rapproche de $A$ et la droite $(AM)$ tend à se confondre avec la tangente à $C_{f}$ en $A$.
\end{minipage}
\hfill
\begin{minipage}{5.6cm}
\begin{center}
\includegraphics{Figures/CoursDerivationFigures.1}
\end{center}
\end{minipage}
%\newpage

\begin{demo}
Soit $T$ la tangente à la courbe au point $A$.
Son coefficient directeur est $f'(a)$ donc l'équation réduite de $T$ est de la forme $y=f'(a)x+b$.

$A\in T$ donc $f(a)=f'(a)a+b$ donc $b=f(a)-af'(a)$.

L'équation de $T$ est donc $y=f'(a)x+f(a)-af'(a)$ soit $y=f'(a)(x-a)+f(a)$.
\end{demo}


\section{Applications de la dérivation}
\subsection{Dérivée et variations}
\begin{prop}
Soit $f$ une fonction dérivable en $a$.
\begin{itemize}
\item Si $f'(a)>0$ Il existe un intervalle $I$ contenant $a$ tel que $f$ est  strictement croissante sur $I$.
\item Si $f'(a)<0$ Il existe un intervalle $I$ contenant $a$ tel que $f$ est  strictement décroissante sur $I$.
\end{itemize}
\end{prop}


\begin{rem}
on utilise souvent les résultats suivants.
\begin{itemize}
\item Si, pour tout $x$ de $I$, $f'(x)>0$ alors $f$ est strictement croissante sur $I$.
\item Si, pour tout $x$ de $I$, $f'(x)<0$ alors $f$ est strictement décroissante sur $I$.
\end{itemize}
\end{rem}

\newpage


\subsection{Extremum local}
\begin{prop}
Soit $f$ une fonction définie sur un intervalle $I$ et $x_{0}\in I$.

On dit que $f(x_{0})$ est un maximum local (respectivement minimum local) de $f$ si l'on peut trouver un intervalle ouvert $J$ inclus dans $I$ et contenant $x_{0}$ tel que pour tout $x\in J$, $f(x)\ie f(x_{0})$ (respectivement $f(x)\se f(x_{0})$).
\end{prop}

\begin{minipage}{7.5cm}
\begin{exemple}
Une fonction est représentée ci-contre.

Son minimum est $-2$

Son maximum est 2.

$-1$ et $-2$ sont des minimums locaux

1 et 2 sont des maximums locaux.
\end{exemple}
\end{minipage}
\hfill
\begin{minipage}{6.5cm}
\begin{center}
\includegraphics{Figures/CoursDerivationFigures.2}
\end{center}
\end{minipage}

\begin{prop}
Soit $f$ une fonction dérivable sur un intervalle $I$ et $x_{0}\in I$.

Si $f(x_{0})$ est un extremum local de $f$ alors $f'(x_{0})=0$.
\end{prop}


\begin{minipage}{7.5cm}
\begin{rem}
 La réciproque de cette propriété est fausse, voir le contre-exemple sur la figure çi contre.
\end{rem}
\end{minipage}
\hfill
\begin{minipage}{6.5cm}
\begin{center}
\includegraphics{Figures/coursderivation3.1}
\end{center}
\end{minipage}

\begin{prop}
Soit $f$ une fonction dérivable sur un intervalle $I$ et $x_{0}\in I$.

Si $f'$ s'annule en changeant de signe en $x_{0}$ alors $f(x_{0})$ est un extremum local de $f$.
\end{prop}

\end{document}
